\scanpage{7}
\uline{Conclusion} La donnée (pour $B$ fixé) d'une catégorie $A$, et d'un couple de foncteurs adjoints $f,g$,
\diagram{\begin{tikzcd}[ampersand replacement=\&] A\arrow[r,shift left,"f"] \& B\arrow[l,shift left,"g"]\end{tikzcd}}
avec $f$ \uline{conservatif} et \deleted{tel que, \unsure{pour les}} \deleted{commutant aux noyaux de doubles \unsure{doubles} flèches} satisfaisant la condition C de 1.12 est équivalente essentiellement à la donnée d'un foncteur $\varphi$ dans $B$, avec les données de Godement $\alpha,\lambda$.

\deleted{De plus, $f$ est pleinement fidèle ssi $\alpha$ est un isomorphisme,}

\uline{Corollaire} $f$ pleinement fidèle \deleted{\unsure{\Longleftrightarrow}}\quad
$\alpha*\varphi\text{ isom}\Longleftrightarrow\lambda\text{ isom}$ [\unsure{alrs|NB} alors $\lambda$ est connu quand on connaît $\alpha$]

\uline{Compléments sur propriétés d'exactitude}
% slanted
\margin{\deleted{Si de plus les \unsure{noyaux}} \uline{Pour} les couples $(f,g)$ \deleted{\unsure{de foncteurs}} adjoints, où \unsure{$\exists$ ds $B$} les Ker \unsure{resp.}, $f$ conservatif et commute aux Ker}

a) Pour tout type de $\varprojlim$ qui sont représentables dans $B$, \deleted{et auxquelles $\varphi$ commute,} les $\varprojlim$ de ce type \deleted{existent} sont repr. dans $A$, et le foncteur $f$ y commute, \inserted{si et seulement si $\varphi$ y commute.}

b) Pour tout type de $\varinjlim$ qui est représentable dans $B$, \unsure{ce} \unsure{\textquotedbl} type est représentable dans $A$, et $f$ y commute (il commute \textquotedbl à \uline{toutes} les $\underline{\varinjlim}$ !)

c) Pour tt propriété d'exactitude de $B$, \deleted{\unsure{faisant|faisant intervenir}} du type commutation de $\varprojlim_J$ \deleted{de type $J$} à $\varinjlim_I$, avec des \deleted{\unsure{typ.}} $\lim$ en question repr. dans $B$, et $\varphi$ comm. aux $\varprojlim_J$, -- la \textquotedbl propriété d'exactitude est vraie dans $A$. \uline{Ex.} Les propriétés d'exactitude a) b) c)

\deleted{d) Supposons que $B$ soit une \unsure{filtration|catégorie} \inserted{du topos} \unsure{continue} $(B_k)$,}
\deleted{\unsure{d)}}
\uline{Théorème} Supposons $B$ un topos. Alors $A$ est un topos et $f$ de la forme $f_0^*$ pour un morphisme de topos $B\longrightarrow A$ (nec. conservatif) sss $\varphi$ \deleted{commute aux \unsure{$\varprojlim$}} \inserted{est exact à gauche \uline{et accessible}}. Pour que $f_0$ soit essentiel, il f. et s. que $\varphi$ commute aux $\varprojlim$ quelconques i.e. aux produits infinis, (ou encore que $\varphi$ ait un adjoint à gauche $\psi$).
% archive sheet number: 6
